\documentclass[a4paper,11pt]{article}
\usepackage{exam-style-341}
\examinehead{341 农业综合知识三 · 数据库技术与应用}{模拟试卷（一）}
\begin{document}
\examcover{341 农业综合知识三}{数据库技术与应用 · 模拟试卷（一）}{数据库技术与应用（50 分）}{50 分}{60 分钟}{本科目只考数据库，与程序设计、网络无关。}

\parttitle{数据库技术与应用（共 50 分）}

% ============================================================
\qtypename{一、单项选择题}{10 题，每题 2 分，共 20 分}
% ============================================================

\q 数据库系统的三级模式结构中，描述数据库中全体数据的全局逻辑结构的是（\quad）
\choicesline{A. 外模式}{B. 模式}{C. 内模式}{D. 存储模式}

\q 数据库系统的三级模式结构中，外模式/模式映像保证了数据的（\quad）
\choicesline{A. 物理独立性}{B. 逻辑独立性}{C. 完整性}{D. 安全性}

\q 数据模型的三要素是（\quad）
\choices{数据结构、数据操作、完整性约束}{实体、属性、联系}{外模式、模式、内模式}{关系、元组、属性}

\q 在关系模式中，能够唯一标识一个元组且不含多余属性的属性组称为（\quad）
\choicesline{A. 主属性}{B. 候选码}{C. 外码}{D. 超码}

\q 在关系数据库的完整性规则中，“若属性 $A$ 是基本关系 $R$ 的主属性，则 $A$ 不能取空值”属于（\quad）
\choicesline{A. 实体完整性}{B. 参照完整性}{C. 用户定义完整性}{D. 域完整性}

\q 设关系模式 $R(A,B,C)$，函数依赖集 $F=\{A\to B,\ B\to C\}$，则 $R$ 的规范化程度最高达到（\quad）
\choicesline{A. 1NF}{B. 2NF}{C. 3NF}{D. BCNF}

\q 设关系 $R(A,B)$ 和 $S(B,C)$，则 $R$ 与 $S$ 进行自然连接后，结果关系中包含的属性个数是（\quad）
\choicesline{A. 2}{B. 3}{C. 4}{D. 5}

\q 事务的 ACID 性质中，“事务一旦提交，它对数据库中数据的改变就是永久性的”指的是（\quad）
\choicesline{A. 原子性}{B. 一致性}{C. 隔离性}{D. 持久性}

\q 在数据库设计中，将 E-R 图转换为关系模式的工作属于（\quad）
\choicesline{A. 需求分析阶段}{B. 概念结构设计阶段}{C. 逻辑结构设计阶段}{D. 物理结构设计阶段}

\q 数据库系统发生事务故障时，对已经发生但尚未提交的事务，恢复子系统的处理是（\quad）
\choices{执行 UNDO，撤销该事务对数据库的所有修改}{执行 REDO，重做该事务对数据库的所有修改}{重装数据库最近的副本}{不做任何处理}

% ============================================================
\qtypename{二、填空题}{5 题，每题 2 分，共 10 分}
% ============================================================

\q 关系代数的五种基本运算是并、差、笛卡尔积、\fillblank{} 和 \fillblank{}。

\q 在关系模型中，二维表中的一行称为 \fillblank{}，一列称为 \fillblank{}。

\q 关系数据库的三类完整性约束中，要求“外码的取值必须是被参照关系中某个元组的主码值，或者为空值”的是 \fillblank{} 完整性。

\q 设关系模式 $R(A,B,C,D)$，函数依赖集 $F=\{AB\to C,\ C\to D,\ D\to A\}$，则属性集 $B$ 关于 $F$ 的闭包 $B^{+}=$ \fillblank{}，属性集 $AB$ 关于 $F$ 的闭包 $(AB)^{+}=$ \fillblank{}。

\q 数据库设计通常分为需求分析、概念结构设计、\fillblank{} 结构设计、\fillblank{} 结构设计、数据库实施、数据库运行与维护六个阶段。

% ============================================================
\qtypename{三、简答题}{2 题，每题 5 分，共 10 分}
% ============================================================

\q 什么是视图？使用视图有哪些优点？
\anslines{3.2cm}

\q 并发操作可能带来哪三类数据不一致性问题？请分别简要说明。
\anslines{3.2cm}

% ============================================================
\qtypename{四、综合题}{2 题，每题 5 分，共 10 分}
% ============================================================

\q 某高校学生选课数据库中三个基本表的结构如下（加下划线的属性为主码）：

\begin{center}
\begin{tabular}{ll}
\toprule
关系名 & 属性（含义） \\
\midrule
Student & \underline{Sno}, Sname, Ssex, Sage, Sdept（学号，姓名，性别，年龄，所在系） \\
Course & \underline{Cno}, Cname, Ccredit（课程号，课程名，学分） \\
SC & \underline{Sno}, \underline{Cno}, Grade（学号，课程号，成绩） \\
\bottomrule
\end{tabular}
\end{center}

请完成下列各题：
\begin{enumerate}[label=（\arabic*）, leftmargin=2.4em, itemsep=1pt]
  \item（1 分）用 \code{CREATE TABLE} 语句定义基本表 \code{SC}，要求定义主码 \code{(Sno, Cno)}，并定义 \code{Sno}、\code{Cno} 分别到 \code{Student}、\code{Course} 的外码。
  \item（2 分）查询选修了“数据库”课程且成绩不低于 90 分的学生的学号、姓名和成绩，结果按成绩降序排列。
  \item（2 分）查询平均成绩高于 80 分、并且至少选修了 2 门课程的学生的学号、选课门数和平均成绩，结果按平均成绩降序排列。
\end{enumerate}
\anslines{5.2cm}

\q 某电商平台的订单管理系统中有一个关系模式

\begin{center}
$R$(订单号, 客户号, 客户名, 商品号, 商品名, 单价, 数量)
\end{center}

\noindent 其语义为：一张订单只属于一个客户，一个客户只有一个客户名；一种商品只有一个商品名和一个单价；一张订单中的一种商品有一个订购数量。相应的函数依赖集为

\begin{center}
$F=\{$\ 订单号 $\to$ 客户号，\ 客户号 $\to$ 客户名，\ 商品号 $\to$ 商品名，\ 商品号 $\to$ 单价，\ (订单号, 商品号) $\to$ 数量 $\}$
\end{center}

请回答下列问题：
\begin{enumerate}[label=（\arabic*）, leftmargin=2.4em, itemsep=1pt]
  \item（1 分）写出 $R$ 的全部候选码，并指出其中的主属性和非主属性。
  \item（2 分）$R$ 最高属于第几范式？说明理由。
  \item（2 分）将 $R$ 分解为 3NF，要求分解既保持函数依赖又具有无损连接性，并说明该分解具有无损连接性的理由。
\end{enumerate}
\anslines{5.6cm}

\end{document}
